\chapter{Pendahuluan}
\label{ch:pendahuluan}

\section{Latar Belakang}
\label{sec:latar-belakang}

Rekam Medis Elektronik (RME) merupakan komponen fundamental dalam penyelenggaraan pelayanan kesehatan di Indonesia. Peraturan Menteri Kesehatan Republik Indonesia (Permenkes) Nomor 24 Tahun 2022 tentang Rekam Medis mewajibkan seluruh fasilitas pelayanan kesehatan (fasyankes) untuk menyelenggarakan rekam medis dalam bentuk elektronik, serta mengharuskan penyelenggaraan tersebut menjamin kerahasiaan, keutuhan (integritas), dan ketersediaan data rekam medis (\cite{permenkes24_2022}). Tinjauan sistematis terhadap implementasi RME turut menunjukkan bahwa keamanan dan kerahasiaan data medis pasien menjadi perhatian utama dalam berbagai studi, mengingat data tersebut bersifat sensitif dan berdampak langsung pada kepercayaan pasien terhadap sistem pelayanan kesehatan (\cite{pramesti2024keamanan}). Kewajiban ini menegaskan bahwa keamanan data RME bukan sekadar kebutuhan teknis, melainkan juga kepatuhan regulasi yang mengikat bagi seluruh penyelenggara pelayanan kesehatan.

Untuk menjawab tantangan tersebut, penelitian sebelumnya mengembangkan DecMed, sebuah kerangka kerja manajemen RME terdesentralisasi yang menggunakan Distributed Ledger Technology (DLT) IOTA sebagai basis kontrol akses (\cite{DecMed}). DecMed mengintegrasikan Capability-Based Access Control (CapBAC), Proxy Re-Encryption (PRE), dan InterPlanetary File System (IPFS) (\cite{benet2014ipfscontentaddressed}) untuk memastikan bahwa setiap akses terhadap data RME hanya dapat dilakukan atas persetujuan aktif pasien, dengan durasi dan cakupan akses yang dapat dibatasi oleh pasien sendiri (\cite{DecMed}).

Meskipun DecMed berhasil menjawab persoalan kontrol akses dan kerahasiaan data, penelitian tersebut secara eksplisit menyatakan bahwa pengelolaan RME juga membutuhkan \textit{audit log} yang andal dan dapat diverifikasi yang mensyaratkan adanya audit log yang mampu merekam seluruh aktivitas sistem secara kronologis, karena pencatatan log merupakan komponen inti dari kontrol audit dan umumnya disyaratkan untuk mendukung kepatuhan regulasi maupun investigasi insiden (\cite{DecMed}). Klaim ini sejalan dengan kerangka kontrol audit pada NIST Special Publication 800-53 Revisi 5, yang menempatkan pencatatan dan peninjauan peristiwa keamanan (kontrol keluarga \textit{Audit and Accountability}) sebagai salah satu persyaratan dasar tata kelola keamanan sistem informasi (\cite{nist80053r5}), sejalan pula dengan kontrol \textit{logging} pada ISO/IEC 27002:2022 (\cite{iso27002_2022}). Catatan audit yang lengkap juga berperan penting sebagai sumber data untuk mendeteksi dan memantau potensi ancaman dalam layanan kesehatan.

Namun demikian, mekanisme pencatatan yang tersedia pada DecMed saat ini belum sepenuhnya memenuhi karakteristik audit log yang baik. Terdapat dua celah utama yang teridentifikasi. \textit{Pertama}, sebagian besar mekanisme yang tercatat pada jaringan IOTA merupakan entri \textit{capability} atau izin akses (\textit{access grant}), yaitu metadata yang menyatakan siapa berhak mengakses data apa dan sampai kapan, bukan catatan atas peristiwa akses yang benar-benar terjadi (\textit{access event}). Akibatnya, sistem belum dapat merekonstruksi riwayat siapa yang benar-benar membaca atau mengubah data RME pada waktu tertentu. \textit{Kedua}, sejumlah operasi pada DecMed, seperti pencabutan akses (\texttt{revoke\_access}) dan pembaruan data administratif maupun klinis (\texttt{update\_administrative\_data}, \texttt{update\_medical\_record}), menyebabkan entri lama pada state tergantikan oleh entri baru alih-alih ditambahkan sebagai riwayat baru yang bersifat \textit{append-only}. Kondisi ini bertentangan dengan karakteristik dasar audit log yang baik, yaitu integritas kronologis dan kelengkapan pencatatan, yang mensyaratkan bahwa setiap catatan peristiwa harus bersifat permanen dan append-only, perubahan kondisi dicatat sebagai entri baru, bukan pembaruan entri lama, sehingga riwayat tindakan dapat diverifikasi keutuhannya (\cite{iso27789_2021}, \cite{nist80053r5}).

Celah tersebut menjadi lebih krusial apabila ditinjau dari perspektif ancaman keamanan sistem. Metode STRIDE (\textit{Spoofing, Tampering, Repudiation, Information Disclosure, Denial of Service, Elevation of Privilege}) merupakan salah satu metode pemodelan ancaman yang umum digunakan untuk mengidentifikasi risiko keamanan pada suatu sistem informasi secara sistematis, yang kemudian dapat dilanjutkan dengan penilaian risiko menggunakan model DREAD (\cite{laksono2021stride}). Salah satu kategori ancaman pada STRIDE, yaitu \textit{repudiation}, merupakan kondisi ketika pengguna sistem dapat menyangkal bahwa ia telah melakukan suatu tindakan tertentu; ancaman ini secara langsung berkaitan dengan ketiadaan atau ketidaklengkapan pencatatan log dan mekanisme audit pada suatu sistem (\cite{laksono2021stride}). Dengan kata lain, kebutuhan pencatatan peristiwa (\textit{event}) pada suatu sistem idealnya tidak ditentukan secara ad hoc, melainkan diturunkan dari analisis ancaman yang relevan terhadap arsitektur sistem yang bersangkutan.

Berdasarkan pemaparan tersebut, Tugas Akhir ini bermaksud menganalisis kebutuhan audit log pada sistem DecMed melalui pendekatan pemodelan ancaman STRIDE, guna menentukan jenis peristiwa (\textit{event}) apa saja yang perlu dicatat agar sistem memiliki audit log yang memenuhi karakteristik integritas dan non-repudiation. Hasil analisis tersebut selanjutnya digunakan sebagai dasar perancangan dan implementasi sebuah layanan audit log (\textit{audit log service}) yang terintegrasi dengan arsitektur DecMed yang telah ada, mengikuti pola layanan pendukung (\textit{supporting service}) yang serupa dengan \textit{proxy re-encryption server} yang sudah digunakan pada DecMed (\cite{DecMed}). Metadata dari setiap peristiwa yang tercatat akan disimpan pada jaringan IOTA untuk menjamin integritas dan keterlacakannya, sedangkan berkas log peristiwa akan disimpan pada IPFS (\cite{benet2014ipfscontentaddressed}), konsisten dengan pola pemisahan penyimpanan on-chain dan off-chain yang telah diterapkan pada DecMed untuk data RME (\cite{DecMed}).

\section{Rumusan Masalah}
\label{sec:rumusan-masalah}

Berdasarkan latar belakang yang telah dipaparkan pada Subbab~\ref{sec:latar-belakang}, terdapat dua rumusan masalah utama dalam Tugas Akhir ini, yaitu

\begin{enumerate}
\item Peristiwa (\textit{event}) apa saja yang perlu dicatat oleh mekanisme audit log pada sistem DecMed?
\item Bagaimana rancangan dan implementasi layanan audit log yang tamper-evident dan sesuai dengan arsitektur DecMed?
\end{enumerate}

\section{Tujuan dan Ukuran Keberhasilan Pencapaian}
\label{sec:tujuan}

Tujuan dari Tugas Akhir ini adalah

\begin{enumerate}
\item Menghasilkan daftar kebutuhan peristiwa (\textit{event requirements}) audit log pada sistem DecMed berdasarkan analisis ancaman menggunakan metode STRIDE.
\item Merancang dan mengimplementasikan layanan audit log yang terintegrasi dengan arsitektur DecMed.
\end{enumerate}

Ukuran keberhasilan pencapaian Tugas Akhir ini didasarkan pada kesesuaian daftar kebutuhan \textit{event} yang dihasilkan terhadap kategori ancaman STRIDE yang teridentifikasi, serta keberhasilan layanan audit log yang dikembangkan dalam mencatat, menyimpan, dan menyajikan kembali peristiwa-peristiwa tersebut secara benar dan tidak dapat diubah (\textit{immutable}).

% CATATAN UNTUK PENULIS:
% Bagian ini masih draf awal dan sebaiknya didiskusikan lebih lanjut
% dengan dosen pembimbing, khususnya terkait metrik/ukuran keberhasilan
% yang lebih terukur (misalnya cakupan skenario pengujian,
% jumlah kategori STRIDE yang tercakup, dsb.)

\section{Batasan Masalah}
\label{sec:batasan-masalah}

Terdapat beberapa batasan masalah dari Tugas Akhir ini, antara lain

\begin{enumerate}
\item Analisis ancaman menggunakan STRIDE dilakukan terhadap arsitektur sistem DecMed yang sudah ada, bukan terhadap perancangan ulang sistem DecMed secara keseluruhan;
\item Layanan audit log yang dikembangkan bersifat pencatat pasif (\textit{passive recorder}) yang menerima dan menyimpan peristiwa, tanpa melakukan validasi otorisasi maupun keputusan kontrol akses terhadap peristiwa yang dicatat; dan
\item Implementasi dilakukan pada lingkungan pengembangan yang sama dengan yang digunakan pada penelitian DecMed sebelumnya.
\end{enumerate}

% CATATAN UNTUK PENULIS:
% Sesuaikan dan tambahkan batasan lain jika diperlukan, misalnya terkait
% cakupan analisis ancaman (apakah mencakup seluruh komponen DecMed atau
% hanya komponen tertentu), maupun batasan pengujian yang akan dilakukan.

\section{Metodologi}
\label{sec:metodologi}

Metodologi yang digunakan dalam Tugas Akhir ini terdiri dari empat tahap pelaksanaan, sebagai berikut.

\begin{enumerate}
\item \textbf{Tahap Perencanaan.} Tahap ini berfokus pada pemahaman arsitektur dan mekanisme sistem DecMed yang telah ada, kajian literatur terkait audit log, threat modeling STRIDE, serta penelitian terkait lainnya sebagai dasar analisis.

\item \textbf{Tahap Analisis Ancaman.} Pada tahap ini dilakukan identifikasi aset, aliran data, dan batas kepercayaan (\textit{trust boundary}) pada arsitektur DecMed, dilanjutkan dengan penerapan metode STRIDE untuk mengidentifikasi ancaman pada setiap komponen. Hasil identifikasi ancaman tersebut kemudian diturunkan menjadi daftar kebutuhan peristiwa (\textit{event requirements}) yang perlu dicatat oleh audit log.

\item \textbf{Tahap Perancangan dan Implementasi.} Berdasarkan daftar
kebutuhan \textit{event} yang telah ditentukan, dilakukan perancangan
arsitektur layanan audit log beserta skema penyimpanan metadata pada
IOTA dan berkas log pada IPFS. Selanjutnya, rancangan tersebut diimplementasikan sebagai layanan yang terintegrasi dengan komponen-komponen DecMed yang sudah ada.

\item \textbf{Tahap Pengujian dan Evaluasi.} Tahap terakhir bertujuan
menguji dan mengevaluasi apakah layanan audit log yang dikembangkanberhasil mencatat seluruh \textit{event} yang telah didefinisikan, serta memvalidasi bahwa data yang tercatat memenuhi karakteristik integritas dan non-repudiation yang diharapkan.
\end{enumerate}

\section{Sistematika Pembahasan}
\label{sec:sistematika}

Berikut merupakan sistematika pembahasan dari Tugas Akhir ini.

Bab I menjelaskan gagasan utama dari Tugas Akhir ini yang meliputi latar belakang, rumusan masalah, tujuan, batasan, metodologi, hingga sistematika pembahasan.

Bab II memaparkan studi literatur yang dilakukan, mencakup RME, audit log, threat modeling STRIDE, kontrol akses, DLT dan IOTA, serta teknologi pendukung lain yang digunakan pada DecMed.

Bab III menjelaskan analisis ancaman menggunakan STRIDE terhadap sistem
DecMed serta rancangan solusi layanan audit log yang diusulkan.

Bab IV menjelaskan hasil implementasi dari rancangan yang dipaparkan pada Bab III, beserta mekanisme dan hasil pengujian yang dilakukan.

Bab V merupakan bab penutup yang berisi kesimpulan dan saran dari penulis.

% PENTING: saya belum membaca full text ketiga paper/dokumen ini secara
% menyeluruh (hanya metadata dan cuplikan hasil pencarian), jadi kamu
% TETAP WAJIB membaca sendiri bagian yang relevan sebelum submit, untuk
% memastikan klaim di teks Bab I ini benar-benar sesuai dengan isi
% aslinya, terutama untuk Al Amin et al. yang detail sitasinya sempat
% saya salah tulis di draf sebelumnya (venue dan judul lengkapnya).
%
% Ringkasan pemakaian tiap key pada Bab I ini:
%   - permenkes24_2022        : kewajiban RME elektronik + jaminan CIA
%   - pramesti2024keamanan    : keamanan & kerahasiaan data RME (tinjauan sistematis)
%   - DecMed                  : arsitektur DecMed, CapBAC/PRE/IPFS, klaim
%                                kebutuhan audit log, pola on-chain/off-chain
%   - benet2014ipfscontentaddressed : peran IPFS sebagai off-chain storage
%   - nist80053r5              : prinsip kontrol audit (keluarga AU),
%                                integritas & non-repudiation log
%   - iso27002_2022            : kontrol logging (8.15) sebagai pelengkap NIST
%   - cert2022insiderthreat    : audit log sbg sumber deteksi/investigasi
%                                penyalahgunaan internal (insider threat)
%   - alamin2025breakglass     : akuntabilitas pada akses darurat (break-glass)
%   - laksono2021stride        : definisi STRIDE+DREAD, kaitan kategori
%                                repudiation dengan kebutuhan logging/audit